\chapter{When the Network Is Not the Calculus}

\section*{11.1}
Loss deletes a transition/message that reliable semantics assumes; it need not create any ill-typed payload.

\section*{11.2}
Examples include deadlock freedom under failures, fairness, fault tolerance, authentication, confidentiality, and application correctness.

\section*{11.3}
PROTO-0 contains no clock state or timeout transition, so a deadline cannot be expressed by its reduction relation.

\section*{11.4}
Use $!\Nat.?\Bool.end!$ against its exact dual and drop the Nat in transit.

\section*{11.5}
Send the right Nat at the right point but interpret it as the wrong account identifier or amount.

\section*{11.6}
Duplication violates an exactly-once/message-multiplicity assumption; reordering violates FIFO/order.

\section*{11.7}
The theorem quantifies over executions generated by its formal transition system. A deployment with loss/crash may not be one of those executions.

\section*{11.8}
Faults can be modeled by extending the state and transitions; they are merely absent from this exact calculus.

\section*{11.9}
Take dual well-typed peers and a lossy transport that drops the first message. The deployment does not complete despite local fidelity, witnessing that fidelity does not imply fault tolerance.

\section*{11.10}
Take a payment trace with a well-typed Nat amount but an application rule requiring a different amount. The protocol accepts the trace; the application invariant is false.

\section*{11.11}
Include delivery (reliable/lossy), ordering, duplication, crash, partition, timing, retry, durability, and adversary assumptions as explicit fields.

\section*{11.12}
Full credit requires comparing concrete failure state/transitions, not merely naming a failure-aware paper.
